%
\documentclass[
  notoc % Suppress Tufte style table of contents.
]{tufte-book}

% Required Tufte packages.
\usepackage{changepage} % or changepage
\usepackage{fancyhdr}
\usepackage{fontenc}
\usepackage{geometry}
\usepackage{hyperref}
\usepackage{natbib}
\usepackage{bibentry}
\usepackage{optparams}
\usepackage{paralist}
\usepackage{placeins}
\usepackage{ragged2e}
\usepackage{setspace}
\usepackage{textcase}
\usepackage{textcomp}
\usepackage{titlesec}
\usepackage{titletoc}
\usepackage{xcolor}
\usepackage{xifthen}

\geometry{paperheight=10in,paperwidth=7in,marginparwidth=30mm,marginparsep=2mm,bindingoffset=10mm,top=10mm,inner=8mm,outer=8mm,bottom=16mm,includehead,includemp}

% Tufte vs. Pandoc workaround.
% Issue: https://github.com/Tufte-LaTeX/tufte-latex/issues/64.
\renewcommand\allcapsspacing[1]{{\addfontfeature{LetterSpace=15}#1}}
\renewcommand\smallcapsspacing[1]{{\addfontfeature{LetterSpace=10}#1}}

% \setmainfont{TeX Gyre Pagella}
\usepackage[utf8]{inputenc}
\usepackage[T1]{fontenc}
\setmainfont{texgyrepagella}[
  Extension = .otf,
  UprightFont = *-regular,
  BoldFont = *-bold,
  ItalicFont = *-italic,
  BoldItalicFont = *-bolditalic,
]

\newfontfamily\JuliaMono{JuliaMono}[
  UprightFont = *-Regular,
  BoldFont = *-Bold
]
\newfontface\JuliaMonoRegular{JuliaMono-Regular}
\newfontface\JuliaMonoBold{JuliaMono-Bold}

\setmonofont{JuliaMono-Medium}[
  Contextuals=Alternate,
  Ligatures=NoCommon
]


\DeclareRobustCommand{\href}[2]{#2\footnote{\url{#1}}}

\usepackage{float}
\floatplacement{figure}{H}

% Listings Julia syntax definition.
\input{/home/runner/.julia/packages/Books/C7CMR/defaults/julia_listings.tex}

% Unicode support.
\input{/home/runner/.julia/packages/Books/C7CMR/defaults/julia_listings_unicode.tex}

% Used by Pandoc.
\providecommand{\tightlist}{%
  \setlength{\itemsep}{0pt}\setlength{\parskip}{0pt}
}
\newcommand{\passthrough}[1]{#1}

\usepackage{longtable}
\usepackage{booktabs}
\usepackage{array}

% Source: Wandmalfarbe/pandoc-latex-template.

\definecolor{linkblue}{HTML}{117af2}
\usepackage{hyperref}
\hypersetup{
  colorlinks,
  citecolor=linkblue,
  linkcolor=linkblue,
  urlcolor=linkblue,
  linktoc=page, % Avoid Table of Contents being nearly completely blue.
  pdftitle={Your Awesome Book},
  pdfauthor={Example author 1; Example author 2},
  pdflang={en-US},
  breaklinks=true,
  pdfcreator={LaTeX via Pandoc}%
}
\urlstyle{same} % disable monospaced font for URLs

\title{Your Awesome Book}
\author{\noindent{Example author 1}\\[3mm] \noindent{Example author
2}\\[3mm] }
\date{}

% Re-enable section numbering which was disabled by tufte.
\setcounter{secnumdepth}{2}

% Fix captions for longtable.
% Thanks to David Carlisle at https://tex.stackexchange.com/a/183344/92217.
\makeatletter
\def\LT@makecaption#1#2#3{%
  \noalign{\smash{\hbox{\kern\textwidth\rlap{\kern\marginparsep
  \parbox[t]{\marginparwidth}{\vspace{12pt}%
\@tufte@caption@font \@tufte@caption@justification \noindent
   #1{#2: }\ignorespaces #3}}}}}}
\makeatother

% Doesn't seem to do anything.
\usepackage{float}
\floatplacement{figure}{H}
\floatplacement{table}{H}

% Reduce large spacing around sections.
\titlespacing*{\chapter}{0pt}{5pt}{20pt}
\titlespacing*{\section}{0pt}{2.5ex plus 1ex minus .2ex}{1.3ex plus .2ex}
\titlespacing*{\subsection}{0pt}{1.75ex plus 1ex minus .2ex}{1.0ex plus.2ex}

\titleformat{\chapter}%
  [hang]% shape
  {\normalfont\huge\itshape}% format applied to label+text
  {\huge\thechapter}% label
  {1em}% horizontal separation between label and title body
  {}% before the title body
  []% after the title body

% Reduce spacing in table of contents.
\usepackage{etoolbox}
\makeatletter
\pretocmd{\chapter}{\addtocontents{toc}{\protect\addvspace{-3\p@}}}{}{}
\pretocmd{\section}{\addtocontents{toc}{\protect\addvspace{-4\p@}}}{}{}
\pretocmd{\subsection}{\addtocontents{toc}{\protect\addvspace{-5\p@}}}{}{}
\makeatother

% Long texts are harder to read than tables.
% Therefore, we can reduce the font size of the table.
\AtBeginEnvironment{longtable}{\footnotesize}

% Some space between paragraphs is necessary because code blocks can output single line paragraphs.
\setlength\parskip{1em plus 0.1em minus 0.2em}

% For justified text.
\usepackage{ragged2e}

% tufte-book disables subsubsections by default.
% Got this definition back via `\show\subsubsection`.

\usepackage{amsfonts}
\usepackage{amssymb}
\usepackage{amsmath}
\usepackage{unicode-math}

% URL line breaks.
\usepackage{xurl}

% Probably doesn't hurt.
\usepackage{marginfix}


\begin{document}

\makeatletter
\thispagestyle{empty}
\vfill
{\Huge\bf
\noindent
\@title
}\\[1in]
{\Large
\noindent
\@author
}
\makeatother

\makeatletter
\newpage
\thispagestyle{empty}
\vfill
{\noindent
\begin{tabular}{l} Example author 1\\ Example institute 1\\ \\ Example author 2\\ Example institute 2\\ \\ \end{tabular}
}
\vfill
{\small
\url{https://booktemplate.huijzer.xyz}

2022-08-06

Creative Commons Attribution-NonCommercial-ShareAlike 4.0 International
}
\makeatother


% Don't remove this or authors will show up in header of every page.
\frontmatter
\mainmatter
\fancyfoot[C]{\url{https://github.com/johndoe/Book.jl}}

\setcounter{tocdepth}{1}
\tableofcontents

% Justify text.
\justifying

% parindent seems to be set from within another class too.
% it is really not useful here because it will also indent lines directly after
% code blocks. Which most of the times not useful.
\setlength{\parindent}{0pt}

\hypertarget{sec:intro}{%
\chapter{Introdução}\label{sec:intro}}

\href{https://julialang.org}{Julia} é uma linguagem relativamente nova
que tem crescido ultimamente, principalmente em áreas relacionadas a
computação científica. Nessa área, Julia compete com Matlab e outros
ambientes parecidos. Mas Julia é uma linguagem genérica que pode ser
utilizada em diferentes aplicações. Nesse sentido, Julia compete com
Python (que, hoje, também compete com Matlab). Um dos aspectos mais
celebrados de Julia é que consegue ter alto desempenho competindo assim
com C/C++/Fortran.

Parafraseando o anuncio da linguagem Julia na web em 2012
(\url{https://julialang.org/blog/2012/02/why-we-created-julia/}), Julia
propõe ser uma linguagem que combina a flexibilidade e dinamismo do
Python, usando uma notação matemática como Matlab, permitindo usar
macros como lisp, processando strings como Perl, tão potente para
álgebra linear quanto Matlab, permitindo interfacear programas como
Shell, sendo fácil de aprender mas não pegando nos nervos de
programadores avançados e sendo interativo e compilado. Jé mencionei que
a linguagem deve ser tão rápida quanto C (ou C++ ou Fortran)?

Passados 10 anos deste anuncio, é possível dizer que, em grande parte,
estes objetivos foram alcançados.

\hypertarget{sec:porquejulia}{%
\section{Por quê Julia?}\label{sec:porquejulia}}

As características acima parecem boas demais para serem verdade. Mesmo
se forem, escolher uma nova linguagem vai muito além das características
positivas da linguagem. Vários outros aspectos são importantes, tanto
para programadores experientes quanto para iniciantes. Aprender uma nova
linguagem não é um investimento trivial e deve ser uma decisão bem
pensada (a não ser que você se divirta aprendendo novas linguagens de
programação\ldots). Assim, algumas coisas devem ser levadas em
consideração:

\begin{itemize}
\tightlist
\item
  Julia apresenta caracterísiticas positivas? Se eu não menti, acho que
  combinar as melhores características de Matlab, Python, C++ me parece
  interessante.
\item
  Julia funciona no meu computador? Estou escrevendo isso no Linux,
  regularmente interfaceio com hardware em Windows e no Raspberry pi,
  muito do desenvolvimento é feito no Mac e dá para rodar no FreeBSD.
  \href{https://julialang.org/downloads/}{Visite e dê uma olhada ná
  página de download!}
\item
  Quanto é que isso vai custar? Nada! Julia é software livre e o código
  está disponível no github \url{https://github.com/JuliaLang/julia}. O
  código usa licensa MIT, bem liberal - dá para fazer praticamente
  qualquer coisa, só precisa mencioanar que está usando código Julia (na
  dúvida, leia a licensa,
  \url{https://github.com/JuliaLang/julia/blob/master/LICENSE.md}).
\item
  Tem ferramentas para resolver o problema XYZ? O ecossistema tem
  crescido bastante. Não chega no nível do Python mas com a tua ajuda
  chegaremos lá! 😉. Na dúvida dê uma olhada
  \url{https://juliapackages.com/} que lista os pacotes
  \emph{registrados}
\item
  Julia tem futuro? Novas versões com melhorias saem regularmente, a
  quantidade de pacotes tem crescido bastante e a velocidade deste
  cresciemento também aumentado. O uso de Julia na indústria também tem
  crescido (\url{https://juliacomputing.com/case-studies/}).
\item
  E tudo o que eu já fiz, tenho que começar do zero? Para pessoas que já
  programam faz algum tempo este é um problema sério. Interfacear com
  C/Fortran usando a biblioteca padrão é bem fácil e direto
  (\url{https://docs.julialang.org/en/v1/manual/calling-c-and-fortran-code/}).
  Se quiser algo mais automático, dê uma olhada em
  \href{https://github.com/analytech-solutions/CBinding.jl}{CBinding} ou
  \href{https://github.com/JuliaInterop/Clang.jl}{Clang}. Existem
  pacotes para interfacear com
  \href{https://github.com/JuliaInterop/CxxWrap.jl}{C++},
  \href{https://github.com/JuliaPy/PyCall.jl}{Python},
  \href{https://github.com/JuliaInterop/RCall.jl}{R} e outras linguagens
  (dê uma olhada em \url{https://github.com/JuliaInterop}).
\item
  E processamento paralelo? Julia permite programação assíncrona e
  paralela usando multi-threading, processamento distribuído e até GPU
  (Julia é particularmente bom nessa área).
\end{itemize}

Tá, você me convenceu. E documentação? A documentação oficial está
localizada em \url{https://docs.julialang.org/}. O manual é uma mistura
de guia de usuário e tutorial. Existem livros e tutoriais que pode podem
ser acessados em \url{https://julialang.org/learning/}. Infelizmente a
maior parte está em inglês e esta é a motivação para este livro.

E se eu precisar de ajuda, o que faço? Em inglês, existem vários canais:
\url{https://julialang.org/community/}. Em particular, vale ressaltar o
fórum, \url{https://discourse.julialang.org/}, que é bem movimentado e
com excelente nível. Realmente recomendo que visite. Infelizmente em
português, as opções são limitadas. Talvez alguém se proponha a fazer
isso?

\hypertarget{certamente-julia-tem-defeitos-tambuxe9m-nuxe9}{%
\section{Certamente Julia tem defeitos também,
né?}\label{certamente-julia-tem-defeitos-tambuxe9m-nuxe9}}

Eu tentei vender uma imagem onde Julia parece ser algo revolucionário.
Acho que não menti mas tenho que avisar que Julia também tem alguns
defeitos. Na minha opinião, existem dois que são percebidos rapidamente.

\hypertarget{first-time-to-plot}{%
\subsection{\texorpdfstring{\emph{First time to
Plot}}{First time to Plot}}\label{first-time-to-plot}}

O primeiro problema é conhecido na comunidade como \emph{First time to
plot}, tempo para o primeiro gráfico. Ou seja, quando o tempo para se
gerar um gráfico \emph{pela primeira vez é longo}. Depois de plotar o
gráfico pela primeira vez, o desempenho é alto.

Na realidade isso ocorre com qualquer código que é executado pela
primeira. Esse tempo excessivo ocorre pois na primeira vez que se
executa esse código, o código precisa ser compilado antes de ser
executado. E essa compilação leva algum tempo. Em algumas aplicações
isso pode ser quase imperceptível. Mas em outras, como quando se plota
um gráfico em um pacote todo escrito em Julia, cada chamada de função
deve ser compilada. Vejamos como exemplo o seguinte código que calcula o
seno de uma sequência numérica e plota o gráfico.

\begin{lstlisting}
julia> using GLMakie

julia> @time begin; x = 0:0.01:1; y = sin.(2π.*x); lines(x,y); end
 31.091122 seconds (82.86 M allocations: 4.430 GiB, 3.00% gc time, 95.87% compilation time)
\end{lstlisting}

31 segundos! Um horror! Mas se este mesmo código for executado
novamente,

\begin{lstlisting}
julia> @time begin; x = 0:0.01:1; y = sin.(2π.*x); lines(x,y); end
  0.014505 seconds (74.12 k allocations: 4.199 MiB)
\end{lstlisting}

o código leva uma fração de segundo, como seria de se esperar em um
computador moderno.

Isso tem melhorado e existem soluções para contornar isso, algumas das
quais serão mostradas neste livro.

\hypertarget{mensagens-de-erro}{%
\subsection{Mensagens de erro}\label{mensagens-de-erro}}

Se por um lado, o \emph{first-time-to-plot} é sentido no primeiro dia,
existe um outro problema que a longo prazo é um pouco mais chato:
mensagens de erro de difícil interpretação. Aqui, novamente, a causa é a
flexibilidade e reúso de código tão comum e incentivado em Julia.
Novamente, isso tem melhorado e com experiência você aprende a
interpretar as mensagens de erro.

\hypertarget{outras-dificuldades}{%
\subsection{Outras dificuldades}\label{outras-dificuldades}}

Cada um tem sua reclamação. As duas mencionadas acima são as minhas.
Outra mencionada com frequência na comunidade é a dificuldade de se
descobrir como fazer algo.

Isso é resultado da documentação que ainda precisa melhorar mas também
das características da linguagem que a torna tão atrativa: reúso de
código extensivo em todo o ecossistema. Em Matlab, quando você abre o
programa, está praticamente tudo disponível e com a documentação ali. Em
Python, se você estiver trabalhando em alguma aplicação científica ou de
ciência de dados, você carrega as bibliotecas numpy, scipy, matplotlib e
pandas e tá tudo lá. Em Julia não existe o equivalente destas grandes
bibliotecas. A mesma funcionalidade está distribuída em um número muito
maior de pacotes. Aí fica mais difícil como fazer algo.

Alguém pode perguntar se isso é necessário ou a melhor alternativa. Pode
parecer que não mas as coisas são mais complicadas. Em alguma aplicação
específica, alguém pode criar um tipo de vetor com características
específicas para esta aplicação. Em Matlab ou Python, seria mais
complicado interfacear estas bibliotecas com a matplotlib por exemplo.
Mas em Julia isso não só é possível (se o criador do pacote tiver
trabalhado de maneira inteligente) como é estimulado.

À medida que livros, tutoriais e documentação específica para diferentes
áreas aparecerem, este problema deve diminuir. Existem propostas para
melhorar isso mas só o tempo dirá se isso será resolvido.

Outras pegadinhas da linguagem serão descritas ao longo deste livro.

\hypertarget{sec:instalauxe7uxe3o}{%
\section{Instalando Julia}\label{sec:instalauxe7uxe3o}}

O jeito mais fácil é simplesmente
\href{https://julialang.org/downloads/}{baixar os binários para a sua
plataforma} e instalar. Você não precisa ser administrador para instalar
e pode usar localmente. Até recomendo isso para iniciantes.

\hypertarget{instalando-no-windows}{%
\subsection{Instalando no Windows}\label{instalando-no-windows}}

Existem algumas opções para a instalação de Julia no Windows. Existem
versões de 32 bits e 64 bits. Em geral, o recomendado instalar a versão
de 64 bits. O suporte é, inclusive melhor. É lógico que existem algumas
circunstâncias onde você precisará usar uma versão de 32 bits - quem
nunca teve que usar uma dll antigona? 32 bits só se você tiver uma
necessidade especial.

Mesmo para Windows 64 bits (x86\_64) existem duas versões: uma versão de
instalação (\emph{installer}) e outra portável (\emph{portable}). A
versão portável é muito útil se você quiser ter Julia em um
\emph{pendrive} e usar em qualquer computador. Já a versão de
instalação, adiciona Julia ao menu iniciar e pode criar ícones na área
de trabalho. Novamente, você pode instalar como usuário normal, o que é
recomendado.

Durante a instalação, o instalador vai perguntar se você quer colocar
Julia nna variável de ambiente \passthrough{\lstinline!PATH!}. Em geral
isso não é recomendado mas para os fins deste livro, é mais simples
adicionar a pasta de executáveis Julia no
\passthrough{\lstinline!PATH!}.

\textbf{Importante} Se você estiver usando Windows 7, será necessário
instalar algo chamado
\href{https://www.microsoft.com/en-us/download/details.aspx?id=54616}{Microsoft
Management Framework}. Sem isso talvez consiga instalar e fazer algo mas
pode ter algumas limitações.

Ao executar Julia pela primeira vez, será criada uma pasta
\passthrough{\lstinline!.julia!} no teu \passthrough{\lstinline!HOME!}.
Nesta pasta estarão instalados os pacotes e suas dependências. Esta
pasta pode crescer bastante. Em um sistema multi-usuário pode ser
interessante compartilhar esta pasta entre diferentes usuários. Isso não
será abordado neste livro.

\hypertarget{instalando-no-linux}{%
\subsection{Instalando no Linux}\label{instalando-no-linux}}

Se você for mais aventureiro, pode compilar o código fonte mas
dificilmente há motivo para isso. Algumas coisas em Julia são
implementadas em C/C++ e muita gente acha que usar um compilador
otimizado para o seu sistema permitirá obter mais um pouquinho de
desempenho. Esqueça! O código Julia \emph{que você escreve} é compilado
usando \href{https://llvm.org/}{LLVM} e usar um outro compilador não vai
melhorar o desempenho. Talvez em alguma plataforma específica seja
necessário compilar do zero (num cluster ou algo parecido) mas este é um
livro introdutório né?

Use os binários disponíveis. Existem binários para processadores intel
com arquitetura x86 (32 bits) e x86\_64 (ou x64 ou AMD64, 64 bits).
Também existem binários para ARM, tanto 64 bits quanto 32 bits: sim, é
possível rodar Julia no Raspberry Pi!

Para linux x86\_64, existe duas possibilidades: glibc e musl. Se você
não sabe o que é isso ou qual a melhor opção, não há dúvida, escolha
glibc!

Após descompactar os binários, você pode colocar em qualquer lugar.
Basta adicionar a pasta \passthrough{\lstinline!julia-1.x.y/bin!} à
variável de ambiente \passthrough{\lstinline!PATH!} (onde
\passthrough{\lstinline!x!} e \passthrough{\lstinline!y!} é a versão
específica que foi baixada). No meu micro, eu costumo copiar esta pasta
para \passthrough{\lstinline!/opt/!} e faço um link simbólico em
\passthrough{\lstinline!/usr/bin/julia!}:

\begin{lstlisting}[language=bash]
sudo ln -s /opt/julia-1.7.3/bin/julia /usr/local/bin/julia
\end{lstlisting}

Se você quer manter tudo local, pode copiar a pasta para
\passthrough{\lstinline!\~/.local/share!} e fazer um link simbólico para
\passthrough{\lstinline!\~/.local/bin!}:

\begin{lstlisting}[language=bash]
cp -r julia-1.7.3 ~/.local/share
ln -s ~/.local/share/julia-1.7.3/bin/julia ~/.local/bin/julia
\end{lstlisting}

\hypertarget{instalando-julia-no-mac}{%
\subsection{Instalando Julia no Mac}\label{instalando-julia-no-mac}}

Não tenho um Mac, mas a internet é tua amiga.

\hypertarget{terminais-melhores}{%
\subsection{Terminais melhores}\label{terminais-melhores}}

O terminal do Windows não é a oitava maravilha do mundo. E no Windows 7,
o terminal é um horror. Neste caso, é interessante usar algo melhor. A
solução mais simples é usar um terminal chamado de Mintty. Uma boa
notícia é que o Git (programa de controle de versões) para windows já
vem com o Mintty: \url{https://gitforwindows.org/}. Após instalar,
procure por Git Bash. O bom desta opção é que se você for usar Julia
mais a sério, você vai ter que usar o Git.

Outra possibilidade é usar o ConEmu \url{https://conemu.github.io/} que
ainda permite usar mais de um terminal em abas.

\hypertarget{fontes}{%
\subsection{Fontes}\label{fontes}}

Com a possibilidade de usar caracteres Unicode, é preciso ter uma fonte
que tenha estes caracteres.
\href{https://github.com/cormullion/juliamono}{A família de fontes
JuliaMono}, criada por membros da comunidade Julia, é bem completa neste
quesito.

\hypertarget{sec:repl}{%
\section{O REPL e Usando Julia como calculadora}\label{sec:repl}}

O ambiente básico de programação Julia é o REPL
(\emph{R}ead-\emph{E}val-\emph{P}rint-\emph{L}oop - laço de leitura,
avaliação e impressão). Basicamente nesse ambiente você digita código e
vê o resultado em tempo real. Este ambiente é extremamente útil ao
desenvolver programas mas é muito bom para se usar como uma calculadora.

Executando Julia, um terminal será aberto. Neste ponto você pode começar
a brincar com Julia. Você pode usar este ambiente como uma excelente
calculadora:

\hypertarget{sec:editores}{%
\section{Editores e ambiente de desenvolvimento}\label{sec:editores}}

O REPL descrito acima é importante e útil mas no final das contas é
impossível (ou quase, não tenho dúvida de que alguém conseguiria me
provar errado\ldots) desenvolver um programa ou aplicativo sem usar um
editor de texto ou ambiente de desenvolvimento. Quando eu falo editor de
texto, \textbf{não} estou falando de Word ou coisa parecida. O ambiente
tem que permitir a edição de textos em formato ASCII ou UTF-8.

Vários editores reconhecem a sintaxe Julia e auxiliam na edição de
arquivos Julia.

No windows o \passthrough{\lstinline!notepad!} pode ser usado não
recomendo de modo algum. Se quiser algo bem simples que usa poucos
recursos mais bastante útil com bastante funcionalidade, pode-se usar o
\href{https://notepad-plus-plus.org/}{Notepad++}. Mesmo se você não for
usar este programa para trabalhar com Julia, é sempre útil.

Naturalmente,
\href{https://github.com/JuliaEditorSupport/julia-emacs}{Emacs} e
\href{https://github.com/JuliaEditorSupport/julia-vim}{Vi} têm suporte
avançado para Julia.

Hoje, talvez o ambiente mais popular e com maior funcionalidade para
Julia seja o \href{https://www.julia-vscode.org/}{Visual Studio Code}. O
visual studio code se vende como uma alternativa moderna e fácil de usar
de ambientes clássicos como Emacs e Vi já citados acima.

Antes do desenvolvimento do editor atom \url{https://atom.io} parar,
este era o ambiente de desenvolvimento Julia mais usado:
\url{https://junolab.org}. Ainda funciona e é uma boa opção para quem já
está familiarizado com o Atom. Para novos usuários, o Visual Studio Code
é recomendado pois é aí que está havendo mais desenvolvimento.

Julia incentiva o uso de caracteres Unicode e no terminal (REPL) tem
\href{https://docs.julialang.org/en/v1/manual/unicode-input/}{comandos
específicos para a entrada de vários caracteres Unicode que lembra
LaTeX}. Um editor que reconheça estes comandos facilita a vida. Visual
Studio Code, Atom, Emacs, Vi, todos, se configurados corretamente,
reconhecem estes comandos.

\hypertarget{sec:notebook}{%
\section{Interfaces tipo Notebook}\label{sec:notebook}}

As interfaces tipo notebook são uma inovação recente em ambiente de
programação. Elas procuram mesclar as melhores características dos
ambientes interativos tipo REPL com os editores mencionados acima.

Dependendo da aplicação, os editores de texto terão que ser usados,
principalmente quando o que se está desenvolvendo é um código. Mas é
muito comum, principalmente em computação científica, o que se quer
``vender'' não é um código específico mas sim uma análise de dados ou
algo parecido. Nesta situação, a abordagem tradicional era escrever
alguns scripts (programas pequenos que são usados poucas vezes) que são
executados a partir do REPL de maneira interativa gerando gráficos e
tabelas que ajudam a analisar os resultados e escrever relatórios.

Neste tipo de aplicação, é muito comum se fazer o que eu chamaria de
computação exploratória. Quando você começa, você não sabe exatamente o
que vai fazer. Pode variar dependendo dos dados de entrada e resultado.
Você calcula algo, faz um gráfico, volta atrás, muda algum parâmetro até
chegar onde quer (quem já não ficou louco procurando o comando no
histórico do REPL?).

Esta é a situação perfeita para o uso de Notebooks! Notebooks misturam
código, documentação (principalmente matemática) e saída (gráficos,
tabelas, etc) de maneira interativa. É uma interface realmente
revolocionária para este tipo de aplicação.

Os Notebooks surgiram com o
\href{https://www.wolfram.com/mathematica/}{Wolfram Mathematica} mas
realmente se tornaram popular com o
\href{https://jupyter.org/}{Jupyter}. Inicialmente desenvolvido para ser
uma interface notebook para python (se chamava IPython notebook
originalmente), membros da comunidade Julia perceberam que este ambiente
podia ser usado com \href{https://github.com/JuliaLang/IJulia.jl}{Julia}
e o nome mudou para refletir esta flexibilidade: Ju(lia) - Py(thon) e R.
Hoje é possível usar os notebooks Jupyter com uma variedade imensa de
linguagens (até \href{https://xeus-cling.readthedocs.io/en/latest/}{C++}
ou \href{https://lfortran.org/}{Fortran}).

\href{https://github.com/fonsp/Pluto.jl}{Pluto} é uma nova interface
notebook para Julia, um pouco diferente do Jupyter. Tem algumas
vantagens e desvantagens, dependendo do gosto, mas é certamente
interessante. Veremos isso mais para frente.

O uso e operação destas interfaces notebook serão descritas em um
capítulo posterios.

\hypertarget{sec:python}{%
\section{Instalando Python}\label{sec:python}}

\emph{Python?} em um livro de Julia??? Python é uma linguagem legal e
tem bastante coisa implementada que pode ser útil. Em particular, a
interface notebook Jupyter é implementada em Python e deverá ser
instalada para que possa ser usada com Julia.

No Linux, em geral Python já está instalado. Basta instalar o Jupyter (e
depois o pacote \passthrough{\lstinline!IJulia!} - veremos isso mais
adiante). Mas este não é o caso no Windows. A solução mais simples é
instalar a
\href{https://www.anaconda.com/products/distribution}{distribuição de
Python Anaconda} que já vem com o Jupyter instalado. O Anaconda é
grande. Então pode ser interessante instalar o
\href{https://docs.conda.io/en/latest/miniconda.html}{miniconda} que é
uma versão do Anaconda minimalista. Aí você instala o que você precisar
diretamente. \href{https://github.com/JuliaPy/PyCall.jl}{O pacote Julia
PyCall} permite que Julia chame código Python diretamente. Em Windows,
este pacote instala o miniconda. Se você quiser usar um Python já
instalado no teu sistema, basta especificar o executável python na
variável de ambiente \passthrough{\lstinline!PYTHON!}. Vale realçar que
o pacote \passthrough{\lstinline!PyCall!} só instalará o miniconda no
Windows. No Linux e outras plataformas usará a versão instalado no
sistema.

\textbf{Usuários do Windows 7:} As versões de Python 3.9 e 3.10 não
rodam em Windows 7! Você vai precisar da versão 3.8 no máximo. Caso você
esteja usando a versão instalada automaticamente com o pacote
\passthrough{\lstinline!PyCall!}, você terá problemas: a instalação não
reclama que o Windows 7 é incompatível com Python 3.9 (ou acima). Mas
sabendo do problema, fica fácil solucionar.

\hypertarget{estrutura-do-livro}{%
\section{Estrutura do livro}\label{estrutura-do-livro}}

O livro está dividido em 3 partes:

\begin{enumerate}
\def\labelenumi{\arabic{enumi}.}
\tightlist
\item
  Usando Julia, onde os fundamentos de programação em Julia são
  explorados
\item
  Resolvendo problemas em Julia onde o ecossistema Julia é explorado, em
  particular a geração de gráficos.
\item
  Julia avançado que mostra recursos mais sofisticados e como gerar
  pacotes.
\end{enumerate}

\hypertarget{sec:primeiro}{%
\chapter{Meu Primeiro Programa em Julia}\label{sec:primeiro}}

Neste capítulo vamos introduzir os conceitos básicos de programação em
Julia. Começaremos usando Julia como uma calculadora, introduziremos
variáveis e seus tipos e vamos ver como se escreve \emph{scripts} em
Julia. Ao final faremos um \emph{tour} da REPL.

\hypertarget{sec:pr-calculadora}{%
\section{Julia como calculadora e o REPL}\label{sec:pr-calculadora}}

Para usar Julia como calculadora, basta clicar no ícone na área de
trabalho (se você escolheu essa opção na instalação) ou selecione Julia
no menu. Para a turma mais \emph{hardcore} (quem gosta de linha de
comando e terminais), digite \passthrough{\lstinline!julia!} (se estiver
no \passthrough{\lstinline!PATH!} se não estiver você vai ter que
digitar o path completo). Você verá algo assim:

\begin{lstlisting}
[pjabardo@durruti ~]$ julia
               _
   _       _ _(_)_     |  Documentation: https://docs.julialang.org
  (_)     | (_) (_)    |
   _ _   _| |_  __ _   |  Type "?" for help, "]?" for Pkg help.
  | | | | | | |/ _` |  |
  | | |_| | | | (_| |  |  Version 1.7.3 (2022-05-06)
 _/ |\__'_|_|_|\__'_|  |  Official https://julialang.org/ release
|__/                   |

julia> 
\end{lstlisting}

Pronto, você já pode usar Julia como uma calculadora:

\begin{lstlisting}
julia> 1+1
2

julia> sqrt(2)
1.4142135623730951

julia> sin(pi/4)
0.7071067811865475
\end{lstlisting}

Esta interface interativa em modo texto é conhecida como \emph{REPL} -
\textbf{R}ead-\textbf{E}val-\textbf{P}rint-\textbf{L}oop, laço de
leitura, interpretação e impressão (se alguém tiver uma tradução menos
horríve, aceito sugestões\ldots). Isso é a descrição básica do que
ocorre nesta interface modo texto:

\begin{enumerate}
\def\labelenumi{\arabic{enumi}.}
\tightlist
\item
  Ler um conjunto de caracteres (\textbf{R}ead)
\item
  Interpretar estes caracteres (\textbf{E}val)
\item
  Imprimir o resultado no terminal (\textbf{P}rint)
\item
  Esperar por mais caracteres para continuar operando (\textbf{L}oop)
\end{enumerate}

No exemplo acima podemos ver que operadores matemáticos são aceitos
(\passthrough{\lstinline!1+1!}) assim como chamadas de funções
(\passthrough{\lstinline!sqrt(2)!}). A maioria das funções matemáticas
comuns e operadores estão disponíveis no REPL. Mas muitas coisas a mais
estão disponíveis na \href{https://docs.julialang.org/en/v1/}{biblioteca
padrão} ou \href{https://juliapackages.com/}{pacotes externos} (veremos
mais sobre isso daqui a pouco).

Os operadores matemáticos mais comuns podem ser vistos a seguir:

\begin{lstlisting}
julia> 2+3
5

julia> 3-2
1

julia> 2*3
6

julia> 4/2
2.0

julia> 2^3
8
\end{lstlisting}

Para tirar a potência de um número, use \passthrough{\lstinline!\^!}.

\hypertarget{sec:pr-tipo}{%
\section{Diferentes tipos de números}\label{sec:pr-tipo}}

No exeamplo anterior, repare na divisão. Existe algo diferente em
relação aos outros exemplos.

\begin{lstlisting}
julia> 2*3
6

julia> 6/3
2.0
\end{lstlisting}

Esse exemplo mostra que números podem ter diferentes tipos. No exemplo
acima, \passthrough{\lstinline!2!} é um número inteiro de 64 bits e
\passthrough{\lstinline!2.0!} é um número de ponto flutuante de 64 bits.
Para se certificar disso, pode-se usar a função
\passthrough{\lstinline!typeof!} para ver o tipo de qualquer objeto
Julia:

\begin{lstlisting}
julia> typeof(3)
Int64

julia> typeof(3.0)
Float64
\end{lstlisting}

\passthrough{\lstinline!Int64!} é um número inteiro com sinal de 64
bits. Já \passthrough{\lstinline!Float64!} é um número de ponto
flutuante de 64 bits.

\hypertarget{nuxfameros-inteiros}{%
\subsection{Números inteiros}\label{nuxfameros-inteiros}}

Pode-se intuir que existem outros tipos de números. Os números inteiros
podem ter diferentes tamanhos, 8, 16, 32, 64, 128 bits ou tamanho
arbitrário. També pode ter sinal.

\hypertarget{sec:pkg}{%
\chapter{Gerenciador de Pacotes}\label{sec:pkg}}

\backmatter

\end{document}
